\section{Averaging as a Probe for Agreement-Driven Structure}
\label{sec:averaging_probe}

Having clarified that our objective is to approximate grammar-driven agreement
structure rather than arbitrary low-rank matrices
(Section~\ref{sec:target_structure}),
we now introduce averaging operations as simple, GPU-friendly probes for isolating
this structure in gradient matrices.

\subsection{Physical Interpretation of Gradient Entries}
\label{sec:physical_gradient}

We briefly recall the physical meaning of the gradient matrix entries in order to
motivate the averaging operations that follow.

Consider the key projection in a transformer layer,
\[
k = W_k x,
\qquad
W_k \in \mathbb{R}^{d\times d},
\]
where $x \in \mathbb{R}^{d}$ is the input hidden state and $k \in \mathbb{R}^{d}$ is
the resulting key vector.
The gradient of the loss with respect to $W_k$,
\[
G = \nabla_{W_k} \mathcal{L} \in \mathbb{R}^{d\times d},
\]
has entries
\[
G_{i\ell} = \frac{\partial \mathcal{L}}{\partial W_{k,i\ell}},
\]
which measure how strongly the training data pushes the connection from input
feature $\ell$ to output feature $i$.

For a batch of tokens, the gradient aggregates per-token outer products,
\[
G = \sum_{j,t} \delta^{(j,t)} \bigl(x^{(j,t)}\bigr)^\top,
\]
where $\delta^{(j,t)}$ is the backpropagated error signal at the projection for
token $(j,t)$.
Thus, $G$ encodes correlations between activation directions and error directions,
aggregated across many tokens and contexts.

\subsection{Column Averaging: Primary Probe of Agreement}
\label{sec:column_averaging}

We now introduce column averaging as the primary operation for probing
agreement-driven structure.

Let $\mathbf{1} \in \mathbb{R}^{d}$ denote the vector of all ones, and define the
column-averaged vector
\[
u \;\triangleq\; G \mathbf{1} \in \mathbb{R}^{d},
\qquad
u_i = \sum_{\ell=1}^{d} G_{i\ell}.
\]

Substituting the token-level gradient expression yields
\[
u
=
\sum_{j,t}
\delta^{(j,t)} \bigl( (x^{(j,t)})^\top \mathbf{1} \bigr).
\]
Importantly, this expression remains a \emph{sum over tokens}.
The multiplication by $\mathbf{1}$ collapses feature-specific input variation into
a scalar weight for each token, while the aggregation over tokens suppresses
semantic variation under the assumptions of
Section~\ref{sec:foundational_assumptions}.

\paragraph{Interpretation.}
Column averaging answers the question:
\begin{quote}
\emph{Which output directions of this layer are being consistently pushed across
many tokens and sequences, regardless of which specific input features caused the
pressure?}
\end{quote}

Because grammatical structure is shared across sequences, grammar-induced gradient
components reinforce coherently in $\nu$.
In contrast, semantic components vary across sequences and cancel in expectation
when tokens are drawn from semantically independent sequences.
As a result, $\nu$ serves as a direct probe of the dominant agreement-driven update
direction.

\paragraph{Why column averaging works.}
The effectiveness of column averaging does not rely on low rank alone.
It relies on the fact that the agreement spike is both \emph{low rank} and \emph{coherent}:
its singular vectors are dense, and its row/column marginals (e.g., $A\mathbf{1}$ and
$A^\top\mathbf{1}$) do not vanish due to cancellation.
In contrast, the residual semantic component is high-rank and comparatively \emph{incoherent}
under batching: token-specific semantic features vary across sequences and cancel in
these first-order aggregated statistics.
This asymmetry is specific to grammar-driven agreement structure and does not hold
for arbitrary low-rank matrices.

\subsection{Row Averaging: Auxiliary Construction}
\label{sec:row_averaging}

For completeness, we define the row-averaged vector
\[
v \;\triangleq\; G^\top \mathbf{1} \in \mathbb{R}^{d},
\qquad
v_\ell = \sum_{i=1}^{d} G_{i\ell}.
\]
In token-level form,
\[
v
=
\sum_{j,t}
x^{(j,t)} \bigl( (\delta^{(j,t)})^\top \mathbf{1} \bigr).
\]

Row averaging aggregates gradient contributions associated with each input feature
across output dimensions.
Empirically and conceptually, this signal is typically weaker and noisier than the
column-averaged signal, as input features mix semantic and contextual information
more heavily.

We emphasize that row averaging is \emph{not} introduced as an independent probe of
agreement structure.
Rather, it plays an auxiliary role in completing a rank-one matrix reconstruction
once the dominant agreement direction has been identified via column averaging.

\subsection{Rank-One Reconstruction from Averages}
\label{sec:rank_one_from_averages}

Given the column-averaged vector $\nu$ and the row-averaged vector $v$, a natural
rank-one approximation of the gradient matrix is
\[
\widehat{G}_1
\;\triangleq\;
\frac{u v^\top}{\mathbf{1}^\top G \mathbf{1}}.
\]

This construction is algebraic rather than structural: it combines the dominant
left- and right-marginal signals into a rank-one matrix.
Its validity rests entirely on the agreement-driven structure described in
Section~\ref{sec:target_structure} and the cancellation assumptions of
Section~\ref{sec:foundational_assumptions}.

In particular, $\widehat{G}_1$ should be interpreted as an approximation to the
agreement component $A$, not as a generic low-rank approximation to $G$.
When these assumptions fail, the reconstruction may be inaccurate, even if $G$ is
low rank.

\paragraph{Summary.}
Averaging serves as a lightweight, GPU-native probe for isolating agreement-driven
structure in gradient matrices.
Column averaging is the primary mechanism by which shared grammatical update
directions are revealed, while row averaging plays a secondary, algebraic role in
constructing a rank-one approximation.
This asymmetric treatment reflects the underlying linguistic and statistical
structure of the data and is essential for the correctness of the proposed method.